% Source: https://artofproblemsolving.com/community/c3h1872147_why_is_this
% Stable identity: aops-topic-1872147-q01
For an integer $m\ge2$ let $s(m)$ be the sum of the prime factors of $m$ with multiplicity, so $f(m)=s(m)$ for composite $m$.

\textbf{Lemma 1: $s(m)\le m$, with $s(m)<m$ for composite $m\ne4$.} For primes $s(m)=m$. If $m=pt$ with $p$ prime and $t\ge2$, then by induction $s(m)=p+s(t)\le p+t$, and $pt-p-t=(p-1)(t-1)-1\ge0$, with equality only for $p=t=2$. So $s(m)\le m$, with equality for composite $m$ only when $m=4$.

\textbf{Lemma 2: $f(n)\ge5$ whenever $n\ge5$.} For prime $n$, $f(n)=n+1\ge6$. A composite $n$ has at least two prime factors, each at least $2$, so $s(n)\ge4$, and $s(n)=4$ only for $n=2\cdot2=4$. So the sequence never leaves $\{5,6,7,\dots\}$.

\textbf{Descent.} Let $n\ge7$.
\begin{itemize}
\item If $n$ is composite, then $f(n)=s(n)<n$ by Lemma 1.
\item If $n$ is prime, then $n$ is odd, $f(n)=n+1$ is even and composite, and
\[
f(f(n))=s(n+1)=2+s\left(\dfrac{n+1}{2}\right)\le2+\dfrac{n+1}{2}<n,
\]
because $n>5$.
\end{itemize}
So whenever the current value is at least $7$, after one or two steps we reach a strictly smaller value, which is still at least $5$ by Lemma 2. A strictly decreasing sequence of positive integers cannot go on forever, so the sequence eventually takes a value in $\{5,6\}$. If it is $6$ we are done, and if it is $5$, the next value is $f(5)=6$. Hence every starting value $n\ge5$ eventually reaches $\boxed{6}$ (after which the values alternate $6,5,6,5,\dots$).
